\section{Distillation Should Specify a Path, Not Only an Endpoint}
\label{sec:motivation}

\subsection{Embedding Distillation as Static State Matching}
\label{sec:static-state-matching}

Let $f_T$ be a teacher embedding model and let a student encoder $f_S$
contain $L$ Transformer layers. For an input sentence $x_i$, pooling the token
states after layer $\ell$ gives a normalized sentence representation
$\mathbf{z}^{(\ell)}_i$. We denote its normalized teacher embedding by
$\boldsymbol{\tau}_i$. The simplest embedding distillation objective constrains
only the terminal representation,
\begin{equation}
    \mathbf{z}^{(L)}_i \longrightarrow \boldsymbol{\tau}_i.
    \label{eq:endpoint-schematic}
\end{equation}
This supervision is efficient and effective, but leaves the internal path of
the student unconstrained. Multi-layer distillation adds targets at a set of
depths $\mathcal{A}$,
\begin{equation}
    \mathcal{L}_{\mathrm{layer}}
    =
    \sum_{\ell\in\mathcal{A}}
    D\!\left(\mathbf{Z}^{(\ell)},\mathbf{T}\right),
    \label{eq:static-layer-matching}
\end{equation}
where rows of $\mathbf{Z}^{(\ell)}$ and $\mathbf{T}$ correspond to examples in
a mini-batch. Relational variants replace or supplement $D$ with discrepancies
between pairwise similarity matrices. These objectives expose the student to
more teacher information, yet each supervised depth remains an independent
state-matching problem.

The common endpoint in Eq.~\eqref{eq:static-layer-matching} contains no notion
of how much semantic transformation is appropriate at depth $\ell$. A weak
constraint may have little effect before the final layers, whereas a strong
constraint can force shallow features to imitate a representation that the
teacher produces only after substantially more computation. Selecting anchor
layers reduces this tension but does not determine the transition between them.
The missing object is a teacher-conditioned rule for local change.

\subsection{Representation Depth as Semantic Time}
\label{sec:semantic-time}

Residual computation suggests indexing representations by a normalized depth
coordinate
\begin{equation}
    t_\ell=\frac{\ell}{L},
    \qquad
    \mathbf{Z}^{(\ell)}\approx\mathbf{Z}(t_\ell),
    \label{eq:semantic-time}
\end{equation}
and viewing consecutive layers as observations along an underlying semantic
trajectory. The corresponding supervision object is a teacher-conditioned
vector field
\begin{equation}
    \frac{d\mathbf{Z}(t)}{dt}
    =F\!\left(\mathbf{Z}(t),\mathbf{T},t\right).
    \label{eq:general-semantic-flow}
\end{equation}
This view changes the question asked at intermediate depths. Static matching
asks whether $\mathbf{Z}^{(\ell)}$ already resembles the terminal teacher
state; trajectory supervision asks whether the observed transition
$\mathbf{Z}^{(\ell)}\rightarrow\mathbf{Z}^{(\ell+1)}$ moves in a direction and
at a pace that reduce teacher discrepancy.

We use ``semantic time'' as an ordered coordinate for supervision, not as an
assumption that a pretrained Transformer is literally the numerical solution of
a latent ODE. The model remains a conventional finite-depth encoder. Continuous
dynamics provide a principled way to construct compatible local targets whose
composition leads toward the teacher endpoint.

\subsection{Connections to Layer-Wise, Geometric, and Flow-Based Distillation}
\label{sec:related-paradigms}

\paragraph{Sentence-embedding distillation.}
Directly regressing teacher sentence embeddings is attractive because teacher
outputs can be cached and reused \citep{reimers2020multilingual}. Contrastive
objectives and multi-stage procedures enrich this signal
\citep{xu2023distillcse,zhang2024jasper}, while recent methods address
architectural and tokenizer mismatch using hidden-state relations or optimal
transport \citep{truong2025emo}. GeoODE-KD retains the efficiency of cached
sentence-level supervision, but uses each teacher output to define both an
endpoint and a direction of evolution.

\paragraph{Layer-wise and geometric transfer.}
Depth-aware embedding methods encourage semantics to emerge at multiple student
depths. ESE trains depth-scalable representations, while TALAS combines upper
layer teacher anchors with adjacent-layer structural propagation
\citep{li2025ese,talas}. These methods motivate our focus on depth but supervise
states or relations at selected locations. GeoODE-KD instead derives every
local transition from one global teacher-conditioned potential. Relational
knowledge therefore affects the direction of representation change, rather
than serving only as an additional state constraint.

\paragraph{Continuous-depth and flow-based transfer.}
Neural ODEs replace a discrete block sequence with a learned differential
equation, and ODE-inspired Transformers use numerical-integration ideas in
architecture design \citep{chen2018neuralode,li2022odetransformer}. Diffusion
and flow-matching approaches have also been used to transport student features
or distributions toward teacher features
\citep{huang2023diffkd,lipman2023flow,shao2024fmkt}. GeoODE-KD does neither: it
learns no velocity network and performs no continuous solve. Its vector field
is the analytic gradient of a teacher-conditioned potential, and the existing
Transformer layers are optimized to agree with a geometric discretization of
that field. This distinction allows the continuous formulation to disappear
entirely at inference time.
